\subsection{Preliminaries}

We first state Vitali's theorem, which lets us convert pointwise convergence to uniform convergence. See~\cite{schiff1993normal} for a reference. 

\begin{fact}[Vitali's convergence theorem]
Let \(h_1,h_2,\ldots\) be holomorphic functions on a connected open set
\(\Omega\). Suppose that:
\begin{enumerate}
    \item the sequence \((h_N)_{N\geq 1}\) is locally bounded on \(\Omega\); and
    \item \(h_N(z)\) converges for every \(z\) in some subset of \(\Omega\)
    having an accumulation point in \(\Omega\).
\end{enumerate}
Then \((h_N)_{N\geq 1}\) converges locally uniformly on \(\Omega\) to a
holomorphic function.
\end{fact}


\subsection{Limiting Krivine rounding schemes}
\label{sec:LimitingKrivine}

In this section, we extend classical Krivine schemes by allowing them to have arbitrary correlations among variables (up to some criteria). We will show that we can state asymptotic statements about classical Krivine schemes using this extended class of schemes, that we call \textit{limiting Krivine schemes}. This gives us analytic machinery for searching and analyzing limits of high dimensional schemes that were previously inaccessible via standard techniques. It is first instructive for us to start with a motivating example. 

\paragraph{Motivating example.}
Consider the two-dimensional classical Krivine rounding scheme
\[
f(X_1,X_2)=\operatorname{sgn}(X_1+\lambda X_2),
\qquad
g(Y_1,Y_2)=\operatorname{sgn}(Y_1+\lambda Y_2),
\]
where \(X_1,X_2,Y_1,Y_2\) are standard Gaussians, \(X_1\) is independent of \(X_2\), \(Y_1\) is independent of \(Y_2\), and
\[
\operatorname{corr}(X_1,Y_1)
=
\operatorname{corr}(X_2,Y_2)
=
t.
\]
In a classical Krivine scheme, every pair of corresponding coordinates has correlation \(t\). Suppose instead that we would like the first pair to have correlation \(t\), but the second pair to have correlation \(t^3\). To obtain this from a classical high-dimensional scheme, define
\[
U_N=\frac{1}{\sqrt N}\sum_{j=1}^N \mathrm{He}_3(X_{2,j}),
\qquad
V_N=\frac{1}{\sqrt N}\sum_{j=1}^N \mathrm{He}_3(Y_{2,j}),
\]
where the pairs \((X_{2,j},Y_{2,j})\) are independent standard Gaussian pairs with correlation \(t\). The Hermite identity gives
\[
\mathbb{E}[U_NV_N]
=
\mathbb{E}[\mathrm{He}_3(X_{2,1})\mathrm{He}_3(Y_{2,1})]
=
t^3.
\]
Moreover, by the central limit theorem, \((U_N,V_N)\) converges in distribution to a standard Gaussian pair \((U,V)\) with correlation \(t^3\). Thus, the classical rounding schemes
\[
F_N=\operatorname{sgn}(X_1+\lambda U_N),
\qquad
G_N=\operatorname{sgn}(Y_1+\lambda V_N)
\]
converge to the fixed-dimensional Gaussian rounding scheme
\[
f(X_1,U)=\operatorname{sgn}(X_1+\lambda U),
\qquad
g(Y_1,V)=\operatorname{sgn}(Y_1+\lambda V),
\]
in which the two coordinate pairs have correlations \(t\) and \(t^3\), respectively.

To use this construction, we need to carefully justify the following analytic steps:
\begin{enumerate}
    \item Pointwise convergence of the correlation functions must be strengthened to locally uniform convergence;
    \item An inverse-majorant $\gamma$-admissibility proved for the limiting scheme must remain valid for all sufficiently large finite-dimensional approximations.
\end{enumerate}

We now generalize this construction and prove both properties. First, we define the notion of allowable correlation maps.  
\begin{definition}[Allowable correlation map]
\label{def:allowable-correlation-map}
A function \(\rho:(-1,1)\to\mathbb{R}\) is an \emph{allowable correlation map}
if it admits an expansion
\[
\rho(t)
=
\sum_{\substack{d\geq 1\\ d\ \mathrm{odd}}} c_d t^d,
\qquad
\sum_{\substack{d\geq 1\\ d\ \mathrm{odd}}}|c_d|
\leq 1.
\]
\end{definition}

We can now state the definition of limiting Krivine schemes. 

\begin{definition}[Limiting Krivine Schemes]
A $k$-dimensional \emph{limiting Krivine scheme} is a tuple $\mathcal{S} = (f, g, \rho)$, where  \(f,g:\mathbb{R}^k\to\{-1,1\}\) are odd measurable functions where the discontinuity sets have Gaussian measure zero, and \(\rho = (\rho_1,\ldots,\rho_k)\) is a tuple of allowable correlation maps.  

The arcsine-normalized correlation function of this scheme is given by 
\[
H_{f, g, \rho}(t)
=
\frac{\pi}{2}\mathbb{E}[f(X)g(Y)].
\]
where \(X,Y\in\mathbb{R}^k\) are jointly Gaussian such that
\[
\mathbb{E}[X_iX_j]
=
\mathbb{E}[Y_iY_j]
=
\delta_{ij},
\qquad
\mathbb{E}[X_iY_j]
=
\delta_{ij}\rho_i(t).
\]
\end{definition}

\begin{remark}
    This is a generalization of the usual high dimensional Krivine schemes with $\rho_i(t) = t$ for all $i$. 

    At this stage, \(H_{f, g, \rho}\) is defined only on \((-1,1)\); in particular, we do not assume that it extends holomorphically to the strip associated with classical Krivine schemes.
\end{remark}

The motivation for this definition is that every limiting Krivine scheme can be approximated arbitrarily well by ordinary finite-dimensional Krivine schemes, as we will show below. This allows us to use the usual projection scheme to round the vectors, and get the inverse-majorant machinery for free. We make this precise in the following theorem.  

\begin{theorem}[Finite-dimensional approximation of limiting schemes]
Let \(\mathcal S=(f,g,\rho_1,\ldots,\rho_k)\) be a limiting Krivine scheme. Then there exists a sequence of ordinary finite-dimensional Krivine schemes
\(
\mathcal S_1,\mathcal S_2,\ldots,
\)
of increasing dimension, whose correlation functions $H_{\mathcal S_1}, H_{\mathcal S_2}, \ldots$ converge to the correlation function $H_{\mathcal S}$ of \(\mathcal S\), uniformly on every compact subset of \((-1,1)\).
\end{theorem}

We dedicate the rest of the section to proving this theorem. First, we prove an equivalence condition on allowability of a function. 

\begin{lemma}[Allowable function criteria]
A map \(\rho\) is allowable if and only if there are odd functions
\(p,q: \R^3 \to \R\), such that
\(\|p\|_2=\|q\|_2=1\) and
\[
\rho(t)=\mathbb{E}[p(U)q(V)],
\]
where \(U,V\in\mathbb{R}^m\) are standard Gaussian vectors with
\(\mathbb{E}[U_iV_j]=t\delta_{ij}\).
\end{lemma}

\begin{proof}
Suppose first that \(p\) and \(q\) are odd and have \(L_2\)-norm one. Write
their Hermite expansions by degree as
\[
p=\sum_{\substack{d\geq 1\\ d\text{ odd}}}p_d,
\qquad
q=\sum_{\substack{d\geq 1\\ d\text{ odd}}}q_d.
\]
Here \(p_d\) and \(q_d\) are the degree-\(d\) Hermite components.

Different Hermite degrees are orthogonal, while degree \(d\) contributes a
factor \(t^d\). Hence
\[
\mathbb{E}[p(U)q(V)]
=
\sum_{\substack{d\geq 1\\ d\text{ odd}}}
\langle p_d,q_d\rangle t^d.
\]
Set \(c_d=\langle p_d,q_d\rangle\). Then
\[
\sum_d|c_d|
\leq
\sum_d\|p_d\|_2\|q_d\|_2
\leq
\left(\sum_d\|p_d\|_2^2\right)^{1/2}
\left(\sum_d\|q_d\|_2^2\right)^{1/2}
=
1.
\]
Thus the resulting correlation map is allowable.

Conversely, suppose that
\[
\rho(t)=\sum_{d\text{ odd}}c_dt^d,
\qquad
s:=\sum_{d\text{ odd}}|c_d|\leq 1.
\]
For \(x=(x_1,x_2,x_3)\in\mathbb{R}^3\), define
\[
p(x)
=
\sum_{d\text{ odd}}
\sqrt{|c_d|}\operatorname{He}_d(x_1)
+
\sqrt{1-s}\,x_2
\]
and
\[
q(x)
=
\sum_{d\text{ odd}}
\operatorname{sgn}(c_d)\sqrt{|c_d|}
\operatorname{He}_d(x_1)
+
\sqrt{1-s}\,x_3.
\]
Both functions are odd. Orthogonality gives \(\|p\|_2=\|q\|_2=1\).

Let \(U,V\in\mathbb{R}^3\) be standard Gaussian vectors with
\(\mathbb{E}[U_iV_j]=t\delta_{ij}\). The terms involving \(U_2\) and \(V_3\)
have zero cross-correlation. Therefore
\[
\mathbb{E}[p(U)q(V)]
=
\sum_{d\text{ odd}}c_dt^d
=
\rho(t).
\]
\end{proof}

We are now ready to prove our theorem, which follows as a direct corollary of the stronger statement we prove below. 

\begin{theorem}[Finite-dimensional approximation of limiting schemes]
Let \((f,g,\rho_1,\ldots,\rho_k)\) be a limiting Krivine scheme.
Then its correlation function extends holomorphically to the unit disk and is
the locally uniform limit there of the correlation functions of ordinary
finite-dimensional Krivine schemes.
\end{theorem}

\begin{proof}
By the above lemma, for each \(i\in\{1,\ldots,k\}\), there exist odd measurable functions 
\(
p_i,q_i:\mathbb{R}^3\to\mathbb{R}
\)
with $\|p\|_2 = \|q\|_2 = 1$ such that, whenever \(U,V\in\mathbb{R}^3\) are joint standard Gaussian vectors with
\(
\mathbb{E}[U_aV_b]=t\delta_{ab},
\)
we have
\[
\mathbb{E}[p_i(U)q_i(V)]
=
\rho_i(t).
\]
We combine these functions into vector-valued maps. Write
\(u\in\mathbb{R}^{3k}\) in \(k\) blocks as
\[
u=\bigl(u^{(1)},\ldots,u^{(k)}\bigr),
\qquad
u^{(i)}\in\mathbb{R}^3,
\]
and define
\[
P(u)
=
\bigl(
p_1(u^{(1)}),\ldots,p_k(u^{(k)})
\bigr),
\]
\[
Q(u)
=
\bigl(
q_1(u^{(1)}),\ldots,q_k(u^{(k)})
\bigr).
\]
For each positive integer \(N\), define functions on
\(\bigl(\mathbb{R}^{3k}\bigr)^N\) by
\[
F_N(u_1,\ldots,u_N)
=
f\left(
\frac{1}{\sqrt{N}}\sum_{r=1}^N P(u_r)
\right)
\]
and
\[
G_N(v_1,\ldots,v_N)
=
g\left(
\frac{1}{\sqrt{N}}\sum_{r=1}^N Q(v_r)
\right).
\]
Since \(P\), \(Q\), \(f\), and \(g\) are odd, the functions \(F_N\) and
\(G_N\) are odd. They therefore define an ordinary finite-dimensional
Krivine scheme.

Let \(H_N\) denote the correlation function of this scheme. Fix
\(t\in(-1,1)\), and let
\[
(U_1,V_1),\ldots,(U_N,V_N)
\]
be independent pairs of standard Gaussian vectors in \(\mathbb{R}^{3k}\)
satisfying
\[
\mathbb{E}[(U_r)_a(V_r)_b]
=
t\delta_{ab}.
\]
Define
\[
A_N
=
\frac{1}{\sqrt{N}}\sum_{r=1}^N P(U_r),
\qquad
B_N
=
\frac{1}{\sqrt{N}}\sum_{r=1}^N Q(V_r).
\]
By the definition of the correlation function,
\[
H_N(t)
=
\frac{\pi}{2}\,
\mathbb{E}\bigl[f(A_N)g(B_N)\bigr].
\]

\begin{claim}
As \(N \rightarrow \infty\), we have \[
(A_N,B_N)
\xrightarrow{d}
(X,Y),
\]
where \((X,Y)\) is jointly Gaussian and satisfies
\[
\mathbb{E}[X_iX_j]
=
\mathbb{E}[Y_iY_j]
=
\delta_{ij},
\qquad
\mathbb{E}[X_iY_j]
=
\delta_{ij}\rho_i(t).
\]
\end{claim}

\begin{claimproof}
The random vectors
\[
\bigl(P(U_r),Q(V_r)\bigr)\in\mathbb{R}^{2k}
\]
are independent, identically distributed, and centered. Because the
\(k\) Gaussian blocks are independent, their covariance matrix is
\[
\begin{pmatrix}
I_k & D(t)\\
D(t) & I_k
\end{pmatrix},
\qquad
D(t)
=
\operatorname{diag}\bigl(
\rho_1(t),\ldots,\rho_k(t)
\bigr).
\]
By the multivariate central limit theorem, we have
\(
(A_N,B_N)
\xrightarrow[N\to\infty]{d}
(X,Y),
\)
where \((X,Y)\) is jointly Gaussian satisfying
\(
\mathbb{E}[X_iX_j]
=
\mathbb{E}[Y_iY_j]
=
\delta_{ij}\) and 
\(
\mathbb{E}[X_iY_j]
=
\delta_{ij}\rho_i(t),
\)
as desired. 
\end{claimproof}

Returning to the proof of the theorem, since the function
\(
(x,y)\longmapsto f(x)g(y)
\)
 is bounded and has discontinuity set with measure 0, it follows that
\[
\mathbb{E}\bigl[f(X_N)g(Y_N)\bigr]
\longrightarrow
\mathbb{E}\bigl[f(X(t))g(Y(t))\bigr].
\]
Consequently,
\[
H_N(t)\longrightarrow H(t)
\]
for every \(t\in(-1,1)\), where
\[
H(t)
=
\frac{\pi}{2}\,
\mathbb{E}\bigl[f(X(t))g(Y(t))\bigr]
\]
is the correlation function of the limiting scheme.

It remains to convert this pointwise convergence to locally uniform
convergence on the unit disk. For each \(N\), the Hermite expansions of
\(F_N\) and \(G_N\) give a holomorphic extension of \(H_N\) to
\[
\mathbb{D}
=
\{z\in\mathbb{C}:|z|<1\}.
\]
More explicitly, if
\[
F_N=\sum_{\alpha}\widehat{F_N}(\alpha)\operatorname{He}_{\alpha},
\qquad
G_N=\sum_{\alpha}\widehat{G_N}(\alpha)\operatorname{He}_{\alpha},
\]
then
\[
H_N(z)
=
\frac{\pi}{2}
\sum_{\alpha}
\widehat{F_N}(\alpha)\widehat{G_N}(\alpha)z^{|\alpha|}.
\]
For \(z\) in the unit disk, the Cauchy--Schwarz inequality gives
\[
\begin{aligned}
|H_N(z)|
&\leq
\frac{\pi}{2}
\sum_{\alpha}
\bigl|
\widehat{F_N}(\alpha)\widehat{G_N}(\alpha)
\bigr|
|z|^{|\alpha|} \\
&\leq
\frac{\pi}{2}
\left(
\sum_{\alpha}
|\widehat{F_N}(\alpha)|^2
\right)^{1/2}
\left(
\sum_{\alpha}
|\widehat{G_N}(\alpha)|^2
\right)^{1/2} \\
&\leq
\frac{\pi}{2}
\|F_N\|_2\|G_N\|_2
=
\frac{\pi}{2}.
\end{aligned}
\]
Here the second inequality uses \(|z|<1\), and the third follows from
the orthonormality of the Hermite polynomials.
Thus, the family \((H_N)_{N\geq 1}\) is locally bounded on \(\mathbb{D}\).
Since it converges pointwise on the interval \((-1,1)\), which has an
accumulation point in \(\mathbb{D}\). Vitali's theorem implies that
\((H_N)\) converges locally uniformly on \(\mathbb{D}\) to a holomorphic
function. On \((-1,1)\), this function agrees with \(H\). It is therefore
the holomorphic extension of the correlation function of the limiting
scheme.
\end{proof}
We now show the following. 
\begin{theorem}[Upper bound from a limiting Krivine scheme]
\label{thm:upper_bound_limiting_scheme}
Let
\(
\mathcal{S}=(f,g,\rho_1,\ldots,\rho_k)
\)
be a limiting Krivine scheme. Let \(H\) be its correlation
function. Suppose that \(H'(0)\neq 0\), so that \(H\) has a local holomorphic inverse
at the origin. Write
\[
H^{-1}(z)
=
\sum_{j=0}^{\infty}a_{2j+1}z^{2j+1},
\]
and define its majorant by
\[
M(r)
=
\sum_{j=0}^{\infty}|a_{2j+1}|r^{2j+1}.
\]
If there exists \(\gamma>0\), within the radius of convergence of
\(H^{-1}\), such that
\(
M(\gamma)=1,
\)
then
\(
K_G
\leq
\frac{\pi}{2\gamma}.
\)
\end{theorem}


The only remaining thing left to show that an inverse-majorant bound for the
limiting correlation function passes to its finite-dimensional
approximations. We first prove this analytic stability statement.

\begin{lemma}[Finite inverse-majorant transfer]
\label{lem:finite-inverse-majorant-transfer}
Let $F_N,F$ be holomorphic on $\D$, suppose that $F_N\to F$ locally
uniformly on $\D$, and assume that $F_N(0)=F(0)=0$. Suppose that $G$ is a
holomorphic inverse branch of $F$ on a neighborhood of
$\overline{D(0,R_0)}$, takes values in $\D$, and satisfies
$
F(G(\zeta))=\zeta,$ with 
$
 |\zeta|\leq R_0.
$
Write
$
G(\zeta)=\sum_{n\geq 1}A_n\zeta^n.
$
If $0<\gamma<R_0$ and
$
\sum_{n\geq 1}|A_n|\gamma^n<1,
$
then, for all sufficiently large $N$, the function $F_N$ has a holomorphic
inverse branch
\[
G_N(\zeta)=F_N^{-1}(\zeta)=\sum_{n\geq 1}A_{n,N}\zeta^n
\]
defined on a neighborhood of $\overline{D(0,\gamma)}$, and
$
\sum_{n\geq 1}|A_{n,N}|\gamma^n<1.
$
Moreover, $A_{n,N}\to A_n$ for every fixed $n$.
\end{lemma}

\begin{proof}
The proof has three steps. First, Rouch\'e's theorem shows that the inverse
branch of $F$ persists for $F_N$. Second, the resulting inverse branches
converge to $G$, which controls every fixed coefficient. Finally, a Cauchy
estimate on a slightly larger disk controls the remaining tail uniformly in
$N$.

Choose radii
\(
\gamma<\rho<r<R<R_0
\)
and set
\(
\Omega=G(D(0,R)).
\)
Because $G$ is an inverse branch of $F$, it is injective and has nonzero
derivative on a neighborhood of $\overline{D(0,R)}$. Hence
\(
\partial\Omega=G(\partial D(0,R)).
\)
For $w\in\partial\Omega$, we therefore have $|F(w)|=R$. Thus, for every
$|\zeta|\leq r$,
\[
|F(w)-\zeta|\geq R-r,
\qquad w\in\partial\Omega.
\]
The contour $\partial\Omega$ is a compact subset of $\D$. Since
$F_N\to F$ locally uniformly, for all sufficiently large $N$,
\[
\sup_{w\in\partial\Omega}|F_N(w)-F(w)|<\frac{R-r}{2}.
\]
Fix such an $N$. For every $|\zeta|\leq r$, Rouch\'e's theorem implies that
$F_N(w)-\zeta$ and $F(w)-\zeta$ have the same number of zeros in $\Omega$,
counted with multiplicity. The equation $F(w)=\zeta$ has exactly one
solution in $\Omega$, namely $w=G(\zeta)$. Therefore $F_N(w)=\zeta$ also has
exactly one solution in $\Omega$.

These solutions define a holomorphic function $G_N$ on $D(0,r)$. One may
write it using the argument principle as
\[
G_N(\zeta)
=
\frac{1}{2\pi i}
\int_{\partial\Omega}
\frac{wF_N'(w)}{F_N(w)-\zeta}\,dw,
\qquad |\zeta|<r.
\]
This is the unique solution of $F_N(w)=\zeta$ in $\Omega$. In particular,
$G_N(0)=0$ and
\(
F_N(G_N(\zeta))=\zeta.
\)
Thus $G_N$ is the inverse branch of $F_N$ at the origin.

The same argument-principle formula holds for $G$. The denominators above
are uniformly bounded away from zero on
$\partial\Omega\times\overline{D(0,r)}$, and local uniform convergence of
$F_N$ implies uniform convergence of $F_N'$ to $F'$ on $\partial\Omega$.
It follows that
\(
G_N\longrightarrow G
\)
uniformly on $\overline{D(0,r)}$. By Cauchy's coefficient formula,
\(
A_{n,N}\longrightarrow A_n
\)
for every fixed $n$.

It remains to control all coefficients at once. Let
\(
B=\sup_{w\in\overline\Omega}|w|.
\)
Since $G_N(D(0,r))\subseteq\Omega$, Cauchy's estimate on the circle
$|\zeta|=\rho$ gives
\(
|A_{n,N}|\leq B\rho^{-n},
\qquad n\geq 1,
\)
for all sufficiently large $N$.

Set
\(
\delta
=
1-\sum_{n\geq 1}|A_n|\gamma^n
>0.
\)
Choose $M$ large enough that
\(
B\sum_{n>M}\left(\frac{\gamma}{\rho}\right)^n
<\frac{\delta}{3}.
\)
Since the first $M$ coefficients converge, for all sufficiently large $N$,
\(
\sum_{n\leq M}|A_{n,N}|\gamma^n
\leq
\sum_{n\leq M}|A_n|\gamma^n
+
\frac{\delta}{3}.
\)
Combining the head and tail bounds gives
\begin{align*}
\sum_{n\geq 1}|A_{n,N}|\gamma^n
&\leq
\sum_{n\leq M}|A_n|\gamma^n
+
\frac{\delta}{3}
+
B\sum_{n>M}\left(\frac{\gamma}{\rho}\right)^n \\
&<
\sum_{n\geq 1}|A_n|\gamma^n
+
\frac{2\delta}{3}
<1.
\end{align*}
This proves the claim.
\end{proof}
We will now prove ~\cref{thm:upper_bound_limiting_scheme}.
% \begin{theorem}[Upper bound from a limiting Krivine scheme]
% \label{thm:upper-bound-from-limiting-scheme}
% Let
% \(
% \mathcal{S}=(f,g,\rho_1,\ldots,\rho_k)
% \)
% be a limiting Krivine scheme. Let $H$ be its correlation function. Suppose
% that $H'(0)\neq 0$, so that $H$ has a local holomorphic inverse at the
% origin. Write
% \[
% H^{-1}(z)
% =
% \sum_{j=0}^{\infty}a_{2j+1}z^{2j+1},
% \]
% and define its majorant by
% \[
% M(r)
% =
% \sum_{j=0}^{\infty}|a_{2j+1}|r^{2j+1}.
% \]
% If there exists $\gamma>0$, within the radius of convergence of $H^{-1}$,
% such that
% \(
% M(\gamma)=1,
% \)
% then
% \(
% \KG
% \leq
% \frac{\pi}{2\gamma}.
% \)
% \end{theorem}

\begin{proof}
We first move from $\gamma$ to an arbitrary smaller radius $\gamma'<\gamma$, where the majorant inequality is strict. We then transfer this strict certificate to a finite-dimensional approximation and finally let $\gamma'$ approach $\gamma$.

Let $H_N$ be the correlation functions of the ordinary finite-dimensional
Krivine schemes supplied by the finite-dimensional approximation theorem.
Thus
\(
H_N\longrightarrow H
\)
locally uniformly on $\D$. Since all of the schemes are odd,
$H_N(0)=H(0)=0$.

The equality $M(\gamma)=1$ need not remain true under a small perturbation,
so we first work at a slightly smaller radius. Fix any
\(
0<\gamma'<\gamma.
\)
Because
\(
a_1=\frac{1}{H'(0)}\neq 0,
\)
the majorant is strictly increasing on the positive real axis. Hence
\(
M(\gamma')<M(\gamma)=1.
\)
Choose $R_0$ such that
\(
\gamma'<R_0<\gamma.
\)
Then $M(R_0)<1$. If $G=H^{-1}$ denotes the inverse power series, then for
$|z|\leq R_0$,
\(
|G(z)|
\leq
M(|z|)
\leq
M(R_0)
<1.
\)
Thus $G$ is holomorphic on a neighborhood of
$\overline{D(0,R_0)}$ and takes values in $\D$. Moreover, the identity
$H(G(z))=z$ holds near the origin. Both sides are holomorphic on
$D(0,R_0)$, so the identity theorem gives
\[
H(G(z))=z,
\qquad |z|<R_0.
\]
Therefore $G$ is an inverse branch of $H$ on this disk.

We may now apply Lemma~\ref{lem:finite-inverse-majorant-transfer} at the
radius $\gamma'$. For all sufficiently large $N$, the inverse branch of
$H_N$ has an expansion
\[
H_N^{-1}(z)=\sum_{n\geq 1}a_{n,N}z^n
\]
satisfying
\[
\sum_{n\geq 1}|a_{n,N}|(\gamma')^n<1.
\]
The classical Krivine criterion applied to this finite-dimensional scheme
therefore yields
\(
\KG\leq \frac{\pi}{2\gamma'}.
\)
This holds for every $\gamma'<\gamma$. Letting $\gamma'$ increase to
$\gamma$ gives
\[
\KG\leq \frac{\pi}{2\gamma},
\]
as claimed.
\end{proof}
